\section*{Chapter \ref{ch:m3:spot-markets} --- Spot Markets}

\begin{solution}{exo:m3:spot-markets:1}
$3{,}000/60{,}000 = 0.05$. Three fees of 10 basis points make a band of $-3\ln(1 - 0.001) \approx 30$ basis
points on each side, plus the spreads: roughly 0.04985 to 0.05015 before spreads.
\end{solution}

\begin{solution}{exo:m3:spot-markets:2}
$60{,}000/60{,}030 = 0.9995$: USDT at about 99.95 cents makes the prices consistent.
\end{solution}

\begin{solution}{exo:m3:spot-markets:3}
Crypto moves freely between venues, so the ether--bitcoin price is kept in line by transfers of coins alone; the
bitcoin price in won against dollars needs the fiat leg, which capital controls and account rules restrict.
\end{solution}

\begin{solution}{exo:m3:spot-markets:4}
$\Phi(-80/(6{,}000\sqrt{30/525{,}600})) = 3.9\%$ at 30 minutes and 10.6\% at 60 minutes.
\end{solution}

\begin{solution}{exo:m3:spot-markets:5}
The benchmark has 1.5 unrounded trades per rounded one. The exchange's authentic trades are about $12\% \times 2.5 =
30\%$ of its trades, so wash trades are about 70\%.
\end{solution}

\begin{solution}{exo:m3:spot-markets:6}
Both legs execute at once, so the gap is locked. The inventory itself remains: its price risk (hedged, usually, on
a futures venue), the failure of the venues holding it, the cost and delay of rebalancing when flows are one-way,
and the capital it ties up.
\end{solution}

\begin{solution}{exo:m3:spot-markets:7}
(Triangular, cross) = (1{,}254, 2{,}953), (794, 1{,}902), (309, 793) and (26, 120). Each snapshot offers 6
triangles (two directions on three venues) but 18 cross-venue routes (three pairs, six ordered venue pairs), and a
cross trade pays two fees against a triangle's three.
\end{solution}

\begin{solution}{exo:m3:spot-markets:8}
Reported volume may be largely \omterm{def:m3:spot-markets:wash}{wash trading}, and even real volume says little about the cost of an order: compare
\omterm{def:m3:spot-markets:depth}{quoted depth} near the mid, measured over time, and the realised slippage of the firm's own trades.
\end{solution}

\begin{solution}{pb:m3:spot-markets:1}
\textbf{1.} $1.001 \times 1.0005$: 1.0015 dollars, 0.15\% of costs. \textbf{2.} $1.65 \times 0.70 \times
\sqrt{60/525{,}600} = 1.23\%$. \textbf{3.} $0.9995 \times 0.995 \times 0.98 = 97.46\%$. \textbf{4.} $1.0015
\times 1.0123/0.9746 - 1 = 4.03\%$. \textbf{5.} $1.15 \times 0.9746/1.0015 - 1 = 11.9\%$ per round trip.
\textbf{6.} 1.95\%. \textbf{7.} 3.28\%. \textbf{8.} About USD~5{,}956. \textbf{9.} Korean exchanges served residents
with local bank accounts; a foreign firm could not open an account or take won out. \textbf{10.} Korean exchanges did
not allow short selling, so the premium could be traded only by bringing coins in. \textbf{11.} It widened: local demand rose
faster than arbitrage capital could flow in through the fiat leg. \textbf{12.} Both legs were crypto, which moved
freely. \textbf{13.} That the constraint was binding for months at a scale no fee could explain: the limit was
permission to move money, not information. \textbf{14.} Only once: selling inventory in Korea delivers won, which
must still leave to buy coins abroad again; the binding constraint is the fiat leg. \textbf{15.} Only if \omterm{def:m3:blockchains-for-traders:stable}{stablecoins}
bought with won could be sold abroad without the same premium; a premium on buying \omterm{def:m3:blockchains-for-traders:stable}{stablecoins} with won is the same
constraint in another form. \textbf{16.} It is the price of a legal restriction, not a failure to notice a gap; it
is an inefficiency only relative to a world without the restriction. \textbf{17.} Residents with capacity under the
rules and access to accounts abroad, who risked price moves in transfer, the legal limits, and the exchanges holding
their funds. \textbf{18.} That prices are made within currency areas, and that the gaps between them measure the cost
of moving money, not the quality of the price. \textbf{19.} \emph{Named result:} 4.03\% with the 2\% remittance cost,
1.95\% without it. \textbf{20.} The time and the permission it takes to move assets, and above all money, between
them.
\end{solution}

\begin{solution}{iq:m3:spot-markets:1}
A cycle through three pairs on one venue that returns more of the starting asset than it spent. It must beat three
taker fees and the spreads, and the venue's market makers keep cross rates inside that band, so what remains is small
and taken by the lowest-fee, fastest participants.
\iqlookfor{the fee band and competition.}
\end{solution}

\begin{solution}{iq:m3:spot-markets:2}
It is a price in USDT, whose dollar value is itself a traded rate that can move, a little normally and a lot in a
\omterm{def:m3:blockchains-for-traders:stable}{depeg}. \omterm{def:m3:blockchains-for-traders:stable}{Stablecoin} balances are marked at market, and a dollar--USDT arbitrage is a position in USDT.
\iqlookfor{the implied cross and marking at market.}
\end{solution}

\begin{solution}{iq:m3:spot-markets:3}
Compare its trade data with regulated venues': first digits of sizes against Benford's law, clustering at round
sizes, the tail of the size distribution; compare volume with \omterm{def:m3:spot-markets:depth}{quoted depth}, price co-movement with the main venues,
and web traffic or on-chain flows; estimate the wash share from rounding against a benchmark.
\iqlookfor{statistical tests plus cross-checks with independent data.}
\end{solution}

\begin{solution}{iq:m3:spot-markets:4}
The cost of closing it may exceed it: fees, slow crediting of deposits or withdrawals, the price risk of the transfer,
limits on moving fiat, account access, or doubts about the dear venue's solvency (a gap can be a credit spread).
\iqlookfor{latency, access and counterparty risk.}
\end{solution}

\begin{solution}{iq:m3:spot-markets:5}
Price risk of the coins (hedge with futures), venue default (limits per venue, sweeping excess), \omterm{def:m3:blockchains-for-traders:stable}{stablecoin} \omterm{def:m3:blockchains-for-traders:stable}{depeg}
(diversify and mark at market), operational risk in transfers (address allow-lists, reconciliation), and rebalancing
cost when flows are one-way (limits on imbalance, a rebalancing budget).
\iqlookfor{market, credit and operational risk each with a limit.}
\end{solution}

\begin{solution}{iq:m3:spot-markets:6}
Estimate, per venue and asset, the flow of opportunities and their one-way drift from history; the cost and time of
moving inventory between each pair; each venue's credit limit. Solve periodically for target inventories that
maximise expected captured edge minus rebalancing and capital costs within the limits (a transport or small linear
program); rebalance when deviations exceed thresholds, preferring cheap chains and quiet hours; monitor fill rates and
refit.
\iqlookfor{forecast flows, costs, limits, and an optimiser with thresholds.}
\end{solution}
